\section{Adaptive Streaming}




\begin{questionbox}{Domanda 64}
Compare QoS and QoE: what they measure and how they correlate. Which network factors impact streaming QoE?
\end{questionbox}


Quality of Service (QoS) describes technical service or network metrics, such as throughput, latency, jitter and packet loss.

Quality of Experience (QoE) describes the user-perceived quality of the multimedia service. It depends on media quality, startup delay, stalls, quality switches, device, context and expectations.

QoS and QoE are correlated but not linearly. For example, packet loss may be invisible if the buffer hides it, while a short throughput drop can severely hurt QoE if it causes rebuffering.



\begin{questionbox}{Domanda 65}
(AdapStrm) Describe the fundamental architectural differences between a “Push-based” streaming system (e.g., RTP/UDP) and a “Pull-based” system (e.g., \textbf{DASH}).
\end{questionbox}


\textbf{Push-based streaming} (RTP/UDP/WebRTC-style):

\begin{itemize}
  \item server sends packets according to its timing,
  \item low latency,
  \item needs real-time transport/control mechanisms,
  \item often used for conferencing/live interactive media.
\end{itemize}

\textbf{Pull-based streaming} (\textbf{DASH}/\textbf{HLS} over HTTP):

\begin{itemize}
  \item client requests segments,
  \item works with HTTP/CDNs/caches,
  \item client selects bitrate based on throughput and buffer,
  \item higher delay but scalable and robust for VoD/streaming.
\end{itemize}



\begin{questionbox}{Domanda 66}
(AdapStrm) Analyze the role of the client-side buffer in the context of stability and Quality of Experience (QoE).
\end{questionbox}


The buffer stores already downloaded media in seconds. It absorbs throughput variation and jitter. Larger buffer improves stability but increases startup delay and latency.

QoE depends on:

\begin{itemize}
  \item initial startup time,
  \item rebuffering events,
  \item duration of stalls,
  \item segment quality,
  \item quality switches.
\end{itemize}



\begin{questionbox}{Domanda 67}
(AdapStrm) Define “Switching Penalty” and discuss its impact on perceived video quality.
\end{questionbox}


Switching penalty is the QoE loss caused by visible changes in quality between consecutive segments. A common score model is:


$$
J(n)=\lambda_1K_n-\lambda_2|K_n-K_{n-1}|-\phi(\Delta_n)
$$


where:

\begin{itemize}
  \item $K_n$ is segment quality,
  \item $|K_n-K_{n-1}|$ penalizes switching,
  \item $\phi(\Delta_n)$ penalizes rebuffering.
\end{itemize}

Frequent quality oscillations can be worse than stable slightly lower quality.

\begin{tcolorbox}[colback=white, colframe=gray!50, height=4.5cm, sharp corners]
\end{tcolorbox}



\begin{questionbox}{Domanda 68}
(AdapStrm) Explain the evolution of the playout buffer level $B(t)$ using a mathematical model. In your explanation, describe the dynamics of the “playback” (draining) phase and the “rebufferization” (stalling) phase.
\end{questionbox}


Buffer level $B(t)$ is measured in playback seconds:


$$
B(t)=L\cdot T_s
$$


where $L$ is number of stored segments and $T_s$ is segment duration.

For coding \textbf{rate} $R_c$ and throughput $S$:


$$
\frac{dB}{dt}=
\begin{cases}
\frac{S}{R_c}-1 & \text{during playback}\\
\frac{S}{R_c} & \text{during rebuffering}
\end{cases}
$$


During playback, the buffer drains at one playback second per real second. During rebuffering, playout stops, so there is no output drain. Stable playback requires, on average:


$$
S>R_c
$$




\begin{questionbox}{Domanda 69}
(AdapStrm) What is the primary motivation for using HTTP-based protocols for video streaming?
\end{questionbox}


HTTP streaming is stateless, CDN-friendly, firewall-friendly and easy to deploy at scale. \textbf{DASH}/\textbf{HLS} can use standard web infrastructure and let the client adapt representation quality.



\begin{questionbox}{Domanda 70}
(AdapStrm) What is the consequence of an \textbf{ABR} algorithm that systematically overestimates the available bandwidth?
\end{questionbox}


If \textbf{ABR} overestimates throughput, it requests segments with too high a bitrate. Download time grows, the buffer empties, and rebuffering/stalling occurs.



\begin{questionbox}{Domanda 71}
(AdapStrm) Which metric is a direct indicator of streaming QoE from the end-user’s perspective?
\end{questionbox}


Direct QoE indicators are startup delay, rebuffering duration/frequency, quality level and quality switches. For users, uninterrupted stable playback is usually more important than low-level network metrics.



\begin{questionbox}{Domanda 72}
(AdapStrm) At what stage in the session lifecycle does a \textbf{DASH} client process the Media Presentation Description (MPD)?
\end{questionbox}


The client processes the MPD at session start, before downloading media segments. The MPD lists periods, adaptation sets, representations, codecs, bitrates, resolutions and segment URLs.



\begin{questionbox}{Domanda 73}
Explain how an \textbf{ABR} (Adaptive Bitrate) algorithm works: \textbf{\textbf{rate}-based} vs \textbf{buffer-based} logic, and the risk of overestimating available bandwidth.
\end{questionbox}


An \textbf{ABR} algorithm is the client-side rule that selects the bitrate representation for the next media segment. Its goal is to maximize QoE by balancing quality, stability and stall avoidance.

\textbf{\textbf{rate}-based} \textbf{ABR} estimates recent throughput and chooses a representation below the available bandwidth, often with a safety margin.

\textbf{buffer-based} \textbf{ABR} uses buffer occupancy: if the buffer is low it chooses conservative bitrates, while if the buffer is high it can request higher quality.

If available bandwidth is overestimated, the client requests segments that are too large. Download time increases, the buffer drains and playback can stall.

---


